% !PRINTED 310
\weberpage{310}
%<<CONT
$x = -\infty$ all successions of signs have been lost; hence in the double sequence of the $f_\nu, F_\nu$, for $x = +\infty$ all $VP$, and for $x = -\infty$ all $PP$.

Accordingly, by \S.\,100, II, III, we can state the following two theorems:

\gesperrt{V. The number of positive roots of $f(x)$ is as great as, or less by an even number than, the number of change-successions in the double sequence (4).}

\gesperrt{VI. The number of negative roots of $f(x)$ is as great as, or less by an even number than, the number of succession-successions in the double sequence (4).}

But the sum of the numbers of the $VP$ and of the $PP$ is equal to the number of successions of signs in the simple sequence
\weq{5}{A_0, A_1, A_2 \ldots A_n,}
and we can therefore combine V. and VI. into the single theorem:

\gesperrt{VII. The number of all real roots of $f(x)$ is as great as, or less by an even number than, the number of successions of signs in the sequence (5).}

The total number of changes of sign and successions of signs in the sequence (5) is as great as the number of all the real and imaginary roots taken together, namely $n$. For the case in which some terms of the sequence (5) vanish is here excluded. Since, moreover, the first and last terms of the sequence (5) have the same sign, the number of changes of sign in this sequence is an even number. We can therefore, without saying anything new, add to VII. the further theorem:

\gesperrt{VIII. The number of imaginary roots of $f(x)$ is at least equal to the number of changes of sign in the sequence (5).}

In the formulation of the last two theorems it is immaterial whether some of the coefficients $a_1, a_2 \ldots a_n$ vanish, provided only that none of the $A_\nu$ is zero.

For the cubic equation
\[
a_0 x^3 + a_1 x^2 + a_2 x + a_3 = 0
\]
we have
\[
A_1 = 2(a_1^2 - 3a_0a_2)
\]
\[
A_2 = 2(a_2^2 - 3a_1a_3),
\]
%CONT>>
